% chapters/06/section-03.tex
\section{Lompatan sebagai Impuls}
\label{sec:lompatan}

Lompatan adalah perubahan kecepatan yang terjadi seketika, bukan gaya yang
bekerja selama beberapa waktu. Bagian ini membahas cara menerapkannya, kapan
boleh diterapkan, dan bagaimana memilih nilainya.

\subsection{Perubahan Kecepatan Seketika}

Melompat berarti memberikan kecepatan awal yang mengarah ke atas, lalu
membiarkan gravitasi memperlambatnya sampai berhenti di titik tertinggi dan
menariknya kembali ke bawah. Karena pada koordinat layar gerak ke atas berarti
pengurangan nilai tegak, kecepatan lompatan diberikan sebagai nilai negatif.

\begin{kodekeluaran}
player_velocity.y = -jump_speed
\end{kodekeluaran}

Kecepatannya \textbf{diganti}, bukan ditambahkan, dan hanya
satu baris. Inilah yang membedakan impuls dari gaya: gaya bekerja terus-menerus
sehingga pengaruhnya ditambahkan setiap bingkai, sedangkan impuls bekerja
sekali sehingga nilainya langsung ditetapkan.

Setelah baris itu dijalankan, gravitasi bekerja seperti biasa pada bingkai
berikutnya, sehingga benda melambat naik, berhenti di puncak, lalu jatuh
kembali. Lintasan yang dihasilkan diperlihatkan pada
Gambar~\ref{fig:gravitasi-lompatan}.

\begin{figure}[htbp]
\centering
% chapters/06/figures/gravitasi-lompatan.tikz
% Lintasan lompatan: kecepatan awal mengarah ke atas, puncak dengan kecepatan
% tegak nol, lalu jatuh dengan kecepatan yang makin besar karena gravitasi.
\begin{tikzpicture}[
  lintasan/.style={very thick, draw=ubPrimary},
  panah/.style={-{Stealth[length=2.4mm]}, very thick, draw=ubAccent},
  panahgrav/.style={-{Stealth[length=2.4mm]}, very thick, draw=ubRed},
  lantai/.style={thick, draw=ubLine},
  label/.style={font=\scriptsize, align=left},
  keterangan/.style={font=\scriptsize\itshape, ubPrimary, align=left}
]

% Lantai.
\draw[lantai] (0,0) -- (8.0,0);

% Lintasan: naik, mencapai puncak, lalu turun.
\draw[lintasan] (0.8,0) .. controls (1.35,1.9) and (1.85,2.6) .. (2.4,2.6);
\draw[lintasan] (2.4,2.6) .. controls (2.95,2.6) and (3.45,1.9) .. (4.0,0);

% Tinggi lompatan.
\draw[ubGreen, dashed] (0.8,0) -- (0.8,2.6);
\draw[ubGreen, dashed] (0.8,2.6) -- (1.55,2.6);
\node[label, ubGreen, anchor=east] at (0.72,1.3)
  {$h = \dfrac{v_0^2}{2g}$};

% Kecepatan awal, mengarah ke atas.
\draw[panah] (0.98,0.3) -- (1.12,1.45);
\node[label, ubAccent, anchor=south east] at (0.95,0.95) {$v_0$};

% Puncak: kecepatan tegak bernilai nol.
\fill[ubPrimary] (2.4,2.6) circle (1.8pt);
\node[label, anchor=south] at (2.4,2.68) {kecepatan tegak nol};

% Kecepatan saat jatuh, mengarah ke bawah dan makin besar.
\draw[panah] (3.88,1.5) -- (3.72,0.35);
\node[label, ubAccent, anchor=north east] at (3.6,1.05) {makin cepat};

% Arah gravitasi yang selalu ke bawah.
\draw[panahgrav] (5.6,2.3) -- (5.6,0.9);
\node[label, ubRed, anchor=west] at (5.75,1.6)
  {gravitasi $g$,\\selalu ke bawah};

% Keterangan arah pada koordinat layar.
\node[keterangan, anchor=north west, text width=52mm] at (5.2,-0.25)
  {Pada koordinat layar, sumbu tegak menuju ke bawah, sehingga gerak ke atas
   bernilai negatif.};
\end{tikzpicture}

\caption{Lintasan lompatan: kecepatan awal, puncak, dan jatuh}
\label{fig:gravitasi-lompatan}
\end{figure}

\subsection{Mengapa Lompatan Dibaca sebagai Peristiwa}

Bagian~4.2 telah membedakan peristiwa dari keadaan tombol. Perbedaan itu di
sini menjadi penentu benar atau salahnya program.

Bayangkan lompatan dibaca memakai \code{pygame.key.get_pressed()}. Selama
tombol lompatan ditahan, keadaannya selalu benar, sehingga pada setiap bingkai
syarat lompatan terpenuhi. Akibatnya benda akan melompat berulang-ulang tanpa
henti selama tombol ditahan, dan pemain dapat terbang ke atas jendela.

Karena itu, lompatan dibaca dari peristiwa \code{KEYDOWN} yang hanya terjadi
sekali pada setiap penekanan.

\begin{kodekeluaran}
elif event.key in (pygame.K_SPACE, pygame.K_UP, pygame.K_w):
    if on_ground:
        player_velocity.y = -jump_speed
        on_ground = False
\end{kodekeluaran}

\begin{catatabox}
Sebagai pedoman yang berlaku umum: \textbf{keadaan tombol} dipakai untuk
tindakan yang berlangsung selama tombol ditahan (berjalan, mengarahkan),
sedangkan \textbf{peristiwa} dipakai untuk tindakan yang terjadi sekali
(melompat, menembak, membuka menu). Membedakan keduanya mencegah dua
kekeliruan yang berlawanan: karakter yang hanya bergerak satu langkah ketika
tombol ditahan, dan karakter yang melompat berkali-kali dari satu penekanan.
\end{catatabox}

\subsection{Syarat Menyentuh Lantai}

Perhatikan syarat \code{if on_ground} pada potongan kode di atas. Tanpa syarat
itu, pemain dapat melompat di udara berkali-kali dan seolah-olah terbang.
Cara mengenali keadaan menyentuh lantai telah dijelaskan pada
Bagian~\ref{sec:gravitasi-lantai}.

Perhatikan pula bahwa variabel \code{on_ground} ditandai salah segera sesudah
lompatan dilakukan. Penandaan itu mencegah lompatan kedua pada bingkai
berikutnya, sebab posisi benda belum berubah dan pemeriksaan lantai mungkin
masih menganggapnya bersandar pada lantai.

\subsection{Memilih Nilai Gravitasi dan Kecepatan Lompatan}

Kedua nilai itu tidak dipilih secara sembarang. Tinggi lompatan ditentukan
oleh keduanya, dan hubungannya dapat dihitung. Pada titik tertinggi, kecepatan
tegak benda bernilai nol. Dengan menggunakan hubungan antara kecepatan dan
tinggi pada gerak berubah beraturan, diperoleh

\begin{equation}
\label{eq:tinggi-lompat}
h = \frac{v_0^2}{2g},
\end{equation}

dengan $h$ tinggi lompatan, $v_0$ kecepatan lompatan, dan $g$ percepatan
gravitasi. Lama benda berada di udara, yaitu dari menjejak lantai sampai
kembali menyentuhnya, adalah

\begin{equation}
\label{eq:lama-di-udara}
t = \frac{2v_0}{g} .
\end{equation}

Kedua rumus itu dapat dibalik untuk keperluan perancangan: apabila Anda
menginginkan tinggi lompatan tertentu, kecepatan lompatannya adalah
$v_0 = \sqrt{2gh}$. Tabel~\ref{tab:rasa-lompatan} memuat tiga pasangan nilai
yang menghasilkan tinggi lompatan yang sama, tetapi dengan rasa yang berbeda.

\begin{table}[htbp]
\centering
\caption{Tiga pasangan nilai gravitasi dan kecepatan lompatan yang menghasilkan tinggi 150 piksel}
\label{tab:rasa-lompatan}
{\setlength{\tabcolsep}{4pt}\footnotesize%
\begin{tabularx}{\textwidth}{@{}C{1.8cm} C{2.2cm} C{2.0cm} C{2.0cm} Y@{}}
\toprule
\headrow \textbf{Gravitasi} & \textbf{Kecepatan lompatan} & \textbf{Tinggi} & \textbf{Lama di udara} & \textbf{Rasa gerak} \\
\midrule
1200 & 600 piksel/detik & 150 piksel & 1,00 detik & Melayang, seperti di bulan. \\
\addlinespace
1800 & 735 piksel/detik & 150 piksel & 0,82 detik & Sedang, banyak dipakai pada permainan platform. \\
\addlinespace
2400 & 849 piksel/detik & 150 piksel & 0,71 detik & Gesit dan tegas, cocok untuk permainan yang cepat. \\
\bottomrule
\end{tabularx}}
\tabsource{Dihitung dengan $v_0 = \sqrt{2gh}$ dan $t = 2v_0/g$ untuk $h = 150$
piksel.}
\end{table}

Tabel tersebut memperlihatkan sesuatu yang penting bagi perancangan: tinggi
lompatan saja tidak menentukan rasa geraknya. Yang menentukan rasa adalah
perbandingan antara gravitasi dan kecepatan lompatan, yaitu lama benda berada
di udara.
